\begin{abstract}
Este trabalho apresenta o \textit{Hephaestus}, um ambiente computacional integrado desenvolvido inteiramente em Python para a modelagem dinâmica simbólica, simulação numérica e comparação de estratégias de controle não-linear aplicadas a manipuladores robóticos seriais e a Sistemas Manipuladores-Veículos Subaquáticos (UVMSs). A \textit{engine} matemática deriva automaticamente, via formulação de Euler-Lagrange com a biblioteca SymPy, as expressões simbólicas fechadas da matriz de inércia $M(q)$, do vetor de Coriolis $C(q,\dot{q})$ e do gradiente de potencial $G(q)$ para cadeias seriais genéricas de $N$ graus de liberdade, sem a restrição da convenção de Denavit-Hartenberg. Para o cenário subaquático, a classe \texttt{RobotMathHydro} estende o modelo com o tensor de massa adicionada $6\times6$ por elo e o potencial de empuxo aparente, em conformidade com a formulação de Fossen. As expressões simbólicas são compiladas em funções NumPy vetorizadas por \textit{lambdify} com Eliminação de Subexpressões Comuns (CSE), viabilizando a avaliação numérica eficiente a cada passo de integração. Dois níveis de paralelismo são implementados: derivação simbólica paralela via \texttt{ProcessPoolExecutor} e avaliação numérica paralela persistente de $M$, $C$ e $G$ durante a simulação. O módulo de simulação resolve a dinâmica direta em malha fechada por um integrador simplético de Euler com \textit{substeps} independentes, incluindo planejamento de trajetórias no espaço cartesiano (segmentos retos e arcos circulares com perfil cúbico), cinemática inversa em malha fechada (CLIK) com Jacobiano $6\times N$ completo para rastreamento de posição e orientação, seis modos de referência de orientação (incluindo SLERP e apontamento para alvo), inicialização robusta por Levenberg-Marquardt com busca em linha adaptativa, tratamento de singularidades por Mínimos Quadrados Amortecidos (DLS) e gestão de redundância por projeção no espaço nulo. Quatro estratégias de controle não-linear são implementadas e comparadas sob interface unificada: Controle por Torque Computado com PID (CTC-PID), Controle por Modo Deslizante com Torque Computado (CT-SMC), algoritmo Super-Twisting (STA) e Controle por Rejeição Ativa de Distúrbios Linear (LADRC) com Observador de Estado Estendido descentralizado. A plataforma conta com interface gráfica em CustomTkinter organizada em três módulos: Modelagem, Simulação e Análise comparativa de sessões. Os experimentos de validação são realizados em um UVMS de 12 graus de liberdade (veículo com 6 GDL + braço antropomórfico com 6 GDL) em ambiente subaquático. Os resultados mostram que todos os controladores atingem rastreamento preciso ao longo de uma trajetória de seis segundos; CT-SMC e LADRC apresentam a maior eficiência energética, com consumo cerca de 35 vezes inferior ao CTC-PID e ao STA; o STA alcança o menor erro quadrático médio, ao custo de maiores picos de torque. O \textit{Hephaestus} é disponibilizado como plataforma didática e de código aberto, preenchendo a lacuna entre a teoria matemática e a simulação de alto desempenho sem dependência de ferramentas proprietárias.

\textbf{Palavras-chave:} Sistemas Manipuladores-Veículos Subaquáticos. Modelagem Simbólica Lagrangiana. Controle Não-Linear. Torque Computado. Modo Deslizante. Super-Twisting. Rejeição Ativa de Distúrbios. Cinemática Inversa em Malha Fechada. Python. SymPy.
\end{abstract}